% New scoped consequence for moving common profiles; no fixed-profile saving.
\section{Why the common profile must stay fixed}\label{prf:core}
The unsigned envelope of Section~\ref{mcs:core} and the unresolved signed
estimate have different profile quantifiers. We give an obstruction to an
$H$-sized bound uniform through the sup norm alone. It uses the \emph{actual}
coefficients and nonprincipal rows, rather than replacing the columns by
arbitrary weights. It does not contradict an estimate for one fixed smooth
profile or one with adequate derivative dependence.

Fix $K=\mathbb Q(\omega)$, $\omega^2+\omega+1=0$, principal $\nu=1$, and a
finite good set $S$ containing the primes above $2,3$. Write
$\zeta=1+\omega$, $\delta=1+2\omega$, so $\zeta$ has order six and
$\delta^2=-3$. Retain the literal raw symbols and normalization
\[
 I_u(W)=D^{-1/2}\sum_{\substack{m\ \mathrm{squarefree}\ (m,S)=1}}
 A_U(m)\chi_m(u)W(Nm/D),\qquad E_u=I_u-P_u.
\]
Here and below squarefree means \emph{ideal} squarefree. For $U^2=D$ the
accepted high-window identity gives
$A_U(m)=\mu(m)(r_U(m)-1)$ on $D<Nm<2D$, with $r_U(m)$ even and
$|A_U(m)|\ge1$. The multiset of prime-ideal norms of a squarefree ideal
is determined by its rational norm. Thus $A_U(m)=A_N$ depends only on
$N=Nm$, including divisor-norm multiplicities. Moreover
\begin{equation}\label{prf:unit}
 \eta(m):=\chi_m(\zeta)=\zeta^{(Nm-1)/6},\qquad (m,S)=1.
\end{equation}
The primitive character $\eta$ has conductor $(2)^2(\delta)^3$, of norm
$432$; its powers $1,5$ have conductor norm $432$, powers $2,4$ norm $27$,
and power $3$ norm $16$. All five powers are nonprincipal. These row
conductors do not remove the raw zero at any prime dividing a sixth factor.

Put $\phi(t)=\exp(1-1/(1-t^2))$ for $|t|<1$ and zero otherwise.
For integer $D$ set $w_N=\operatorname{sgn}(A_N)\overline{\zeta^{(N-1)/6}}$
when a good squarefree ideal of norm $N$ exists and $D<N<2D$, and zero
otherwise. Define one common, complex profile
\begin{equation}\label{prf:profile}
 W_D(y)=\sum_{D<N<2D}w_N\phi(4(Dy-N)).
\end{equation}
Its disjoint supports lie strictly inside $(1,2)$, its sup norm is one for
large $D$, and $\|W_D^{(j)}\|_\infty\ll_jD^j$. In particular the endpoints
have value zero. This is a profile varying with scale, independent of the
row, not a separate choice for each ideal or each row.

\begin{proposition}[A moving-profile lower bound]\label{prf:lower}
For every fixed $0<\theta<1$, take $H=D^\theta$, $U=\sqrt D$, and
$X=H^{1/6}$. As integer $D\to\infty$,
\begin{equation}\label{prf:lower-bound}
 \sum_{0<Nu\le H}|E_u(W_D)|^2\ge(c_S+o(1))DH^{1/6},\qquad c_S>0.
\end{equation}
This lower bound already comes from literal rows
$u_v=\zeta\beta_v^6$, where $v$ is an arbitrary good ideal of norm at most
$X$ and $\beta_v$ is a generator. These rows have nonprincipal primitive
character $\eta$, so $P_{u_v}=0$. Their exact values are
\begin{equation}\label{prf:row}
 E_{u_v}(W_D)=D^{-1/2}
 \sum_{\substack{m\ \mathrm{squarefree},\ (m,Sv)=1\\D<Nm<2D}}
                      |A_U(m)|.
\end{equation}
Every collision with $v$ remains deleted, even when the primitive conductor
drops to $432$. No squarefree condition is imposed on $v$.
\end{proposition}
At $H=D^{3/4}$ the lower exponent is $9/8$, whereas the current upper
exponent is $7/6$. Thus this does not establish sharpness of $7/6$ or a new
upper saving. For $0<\theta<2/3$, the term $DH^{1/6}$ dominates the current
unsigned envelope; it is necessary up to subpower losses in this
\emph{varying-profile} class, even on nonprincipal rows.

For $k\ge0$ put
$M_k(W)=\sum_{j=0}^k\|(y\partial_y)^jW\|_\infty$.
If, for every $\varepsilon>0$, a bound uniform over all common profiles
with constants independent of $W,D$ were to assert
\[
 \sum_{0<Nu\le H}|E_u(W)|^2
 \ll_\varepsilon M_k(W)^\alpha(HD)^\varepsilon H,
 \qquad \alpha\ge0,
\]
then necessarily $k\alpha\ge1-5\theta/6$. At $\theta=3/4$ this is
$k\alpha\ge3/8$. For a single fixed $\varepsilon$ the necessary condition
is only $k\alpha+(1+\theta)\varepsilon\ge1-5\theta/6$.
These are necessary cost comparisons, not sufficient estimates. A
fixed-$W$ result, a class with bounded seminorms, or the existing
$M_4(W)^2$ dependence is not refuted.

The same rows expose a concrete centering error in the joint completion:
removing the orientation mean of the actual correlation can delete a
nonzero signed contribution on \emph{nonprincipal} rows. The extended proof
retains the literal unit phase, gives an exact four-orientation example,
and explains why orientation projection alone supplies no saving. The
standalone \href{../profile-uniformity.html}{profile and coherent-mean inspector}
shows the complete finite norm groups, arbitrary sixth-factor masks and
exact residual quantities. Its finite tables are diagnostics of this
construction, not asymptotic estimates, training or GRH evidence.
